\section{土地管理法}
\begin{enumerate}
  \item 内容非常庞杂。
  \item 国务院，授权的地方政府 \hz 地方政府职能部门行使国有土地的所有权。
  \item 集体。村小组、村、乡镇。
  \item 农用地的承包。结合民法。
  \item 国土空间规划。了解即可。
  \item 耕地保护。
        \begin{enumerate}
          \item 占用耕地的补偿制度。用多少开多少；自己开 \hz 耕地开垦费。
          \item 永久基本农田。
                \begin{enumerate}
                  \item 80\%；以乡 \hz 镇为单位划定；具体到地块。
                  \item 改变用途 \hz 征收 \yjt 找国务院。
                  \item 禁止挖塘养育，种果树。
                  \item 禁止闲置、荒芜。1 年 \yjt 能种就种；1 年+ \yjt 闲置费；\zd{2 年 \yjt 无偿收回}。
                \end{enumerate}
        \end{enumerate}
  \item （国有）建设用地。
        \begin{enumerate}
          \item 转用审批。改变用途。
                \begin{enumerate}
                  \item 变永久基本农田 \yjt 国务院。
                  \item 非\blh。\yi 规划之内转用，规划谁批找谁；\er 之外，国务院 \hz 政府。
                \end{enumerate}
        \end{enumerate}
  \item 土地征收。改变的是所有权。公共利益。权限：
        \begin{enumerate}
          \item 国务院 \yjt 耕地超过 35 公顷；其他土地超过 70 公顷。省级政府 \yjt 排除法。
          \item 征收农用地的，先办农用地转用审批。
          \item 听证会。多数认为不合法。
          \item 补偿。土地补偿费、安置补助费、住宅、地上附着物、青苗补偿费、\zd{社会保障费用}。
          \item 补偿标准。省级指定，区片综合地价；3 年调整。
        \end{enumerate}
  \item 临时用地。县以上政府自然资源部门批准；城市规划区的 \yjt 城乡规划部门；不超 2 年。
  \item 土地使用权的收回。不管它。
  \item 宅基地。分布给内部成员。
        \begin{enumerate}
          \item 一户一宅；户有所居，多户并居。
          \item 不能占用永久基本农田。
          \item 乡镇政府审批。
          \item 出卖、出租、赠予住宅之后，再申请宅基地的不予批准。
          \item 自愿有偿退出。
        \end{enumerate}
  \item 集体建设用地入市。登记为集体经营性建设用地才行。
        \begin{enumerate}
          \item 村民会议决定。2/3 绝对多数同意。
          \item 再流转。
        \end{enumerate}
  \item 土地纠纷及解决途径。考得少。
        \begin{enumerate}
          \item 确权 \yjt 县协商 \hz 行政诉讼。
          \item 侵权 \yjt 自然资源局调处 \hz 民事诉讼。
          \item 承包 \yjt 农村土地承包仲裁机构。
          \item 行政纠纷。行政法相关。
        \end{enumerate}
\end{enumerate}
